\chapter{Hipótesis y objetivos} % Main chapter title
\label{sec:Hipotesis} % For referencing the chapter elsewhere, use


 
 
Si la función de costo del motor de optimización de SimpliRoute se ajusta con un vector de preferencias inferido de forma individual para cada conductor, a partir de la discrepancia entre las rutas planificadas y las trazas GPS históricas, entonces la adherencia a la secuencia de paradas planificada aumenta respecto de la planificación actual, cuya tasa de referencia es de un 14,19\%.

 
\section{Objetivo general}
 
Desarrollar un marco metodológico que infiera las preferencias individuales de los conductores de SimpliRoute, a partir de la discrepancia entre las rutas planificadas y las trazas GPS históricas, que distinga el efecto estructural de la asignación de zona del efecto personal, y que integre dichas preferencias en el motor de optimización de ruteo con el fin de incrementar la adherencia a la secuencia de paradas planificada.
 
\section{Objetivos específicos}
 
\begin{enumerate}
    \item Caracterizar y cuantificar la discrepancia entre las rutas planificadas por SimpliRoute y las trazas GPS ejecutadas por los conductores, estableciendo la línea base de adherencia y las métricas de evaluación.
    \item Formular un modelo que infiera un perfil de preferencias a nivel individual para cada conductor a partir de dicha discrepancia.
    \item Incorporar al modelo un mecanismo que separe el efecto estructural de la zona del efecto personal del conductor.
    \item Integrar las preferencias inferidas en el motor de optimización, ya sea como una reponderación de la función de costo o como una restricción blanda aprendida, sin requerir el rediseño del sistema de ruteo.
    \item Evaluar el incremento de adherencia del enfoque propuesto respecto de la planificación actual y respecto de un modelo de preferencias agregadas, utilizando las métricas definidas en el primer objetivo.
\end{enumerate}
 